\subsection{完备空间的积与逆极限}

命题 10 完备一致空间的每个积都是完备的。反之，若非空一致空间的一个积是完备的，则每个因子都是完备一致空间。

第一个断言是积空间上 Cauchy 滤子与收敛滤子的刻画（第 1 目命题 4 推论 2，以及第一章 § 7 第 6 目命题 10 推论 1）的推论。反之，设

\[
\mathbf {X} = \prod_ {i \in I} \mathbf {X} _ {i}
\]

是完备的（诸 \(X_{t}\) 非空），并设 \(F_{x}\) 是 \(X_{x}\) 上的一个 Cauchy 滤子。对每个 \(t \neq x\)，设 \(F_{t}\) 是 \(X_{t}\) 上的一个 Cauchy 滤子，并考虑积滤子（第一章 § 6 第 7 目）

\[
\mathfrak {F} = \prod_ {i \in I} \mathfrak {F} _ {i} \text { on } X;
\]

\(\mathfrak{F}\) 是 Cauchy 滤子（第 1 目命题 4 推论 2），因而是收敛的，从而每个 \(\mathrm{pr}_{\mathbf{x}}\mathfrak{F} = \mathfrak{F}_{\mathbf{x}}\) 也收敛（第一章 § 7 第 6 目命题 10 推论 1）。

推论 设 \((\mathbf{X}_{\alpha}, f_{\alpha\beta})\) 是一致空间的逆系统。若诸 \(X_{\alpha}\) 是豪斯多夫的且完备的，则 \(X = \varprojlim X_{\alpha}\) 也是。

因为 X 是豪斯多夫的，并且可以与 \(\prod_{i\in I}X_{\alpha}\) 的一个闭子空间恒同

（第一章 § 8 第 2 目命题 7 推论 2）；于是该推论由命题 10 与命题 8（第 4 目）得出。

完备豪斯多夫一致空间 \(X_{\alpha}\) 的逆极限可以是空的，即使所有 \(X_{\alpha}\) 都非空且所有 \(f_{\alpha\beta}\) 都是满射，离散空间的情形即表明这一点（《集合论》，第三章，§ 1，习题 32）。然而，我们有下述定理：

定理 1（Mittag-Leffler） 设 \((\mathbf{X}_{\alpha}, f_{\alpha\beta})\) 是完备豪斯多夫一致空间的逆系统，其指标集 I 是具有可数共尾子集的有向集；再设对每个 \(\alpha \in I\)，\(X_{\alpha}\) 具有可数的一致邻域基本系 (*)。最后，设对每个 \(\alpha \in I\)，存在满足下述条件的指标 \(\beta \geqslant \alpha\)：

\((\mathbf{ML}_{\alpha \beta})\) 对每个 \(\gamma \geqslant \beta, f_{\alpha \gamma}(\mathbf{X}_{\gamma})\) 在 \(f_{\alpha \beta}(\mathbf{X}_{\beta})\) 中稠密。

设 \(\mathbf{X} = \varprojlim \mathbf{X}_{\alpha}\)，并设 \(f_{\alpha}\) 是典范映射 \(\mathbf{X} \to \mathbf{X}_{\alpha}\)。则对每个 \(\alpha \in \overleftarrow{\mathbf{I}}\) 与每个 \(\beta \geqslant \alpha\)（满足 \((\mathbf{ML}_{\alpha\beta})\)），\(f_{\alpha}(\mathbf{X})\) 在 \(f_{\alpha\beta}(\mathbf{X}_{\beta})\) 中稠密（因而 \(\mathbf{X}\) 非空，只要诸 \(\mathbf{X}_{\alpha}\) 均非空）。

设 \((\lambda_{n})\) 是 I 中共尾的指标序列。从指标 \(\alpha_0 \in \mathbf{I}\) 出发，递归地定义递增序列 \((\alpha_{n})\)，使 \(\alpha_{n} \geqslant \lambda_{n}\) 且 \((\mathrm{ML}_{\alpha_{n}, \alpha_{n+1}})\) 为真。显然序列 \((\alpha_{n})\) 在 I 中共尾。我们将用 \(f_{mn}\) 代替 \(f_{\alpha_m\alpha_n}\)（当 \(m \leqslant n\) 时），并令 \(f_{n,n+1}(\mathbf{X}_{\alpha_{n+1}}) = \mathbf{Y}_n\)。于是，若 \(m \leqslant n\)，则 \(f_{mn}(\mathbf{Y}_n)\) 在 \(\mathbf{Y}_m\) 中稠密；因为按定义，\(f_{m,n+1}(\mathbf{X}_{\alpha_{n+1}})\) 在

\[
f _ {m, m + 1} \left(\mathrm{X} _ {\alpha_ {m + 1}}\right) = \mathrm{Y} _ {m}, \quad \text { and } \quad f _ {m, n + 1} \left(\mathrm{X} _ {\alpha_ {n + 1}}\right) = f _ {m n} \left[ f _ {n, n + 1} \left(\mathrm{X} _ {\alpha_ {n + 1}}\right) \right] = f _ {m n} \left(\mathrm{Y} _ {n}\right).
\]

对 \(n\) 与 \(k\) 用归纳法，可以定义闭对称一致邻域的基本系 \((\mathbf{V}_{kn})_{k\in \mathbb{N}}\)（\(\mathbf{X}_{\alpha_n}\) 的，对每个 \(n\)），使得

\begin{enumerate}
\def\labelenumi{(\arabic{enumi})}
\setcounter{enumi}{1}
\tightlist
\item
\end{enumerate}

\[
\begin{array}{c} \mathbf {V} _ {k + 1, n} ^ {2} \subset \mathbf {V} _ {k n}, \\ (f _ {n, n + 1} \times f _ {n, n + 1}) (\mathbf {V} _ {k, n + 1}) \subset \mathbf {V} _ {k n}. \end{array}\tag{3}
\]

实际上，设 \((\mathbf{U}_{kn})_{k\in\mathbb{N}}\) 是 \(X_{\alpha_{n}}\) 的一个一致邻域基本系。若假设对给定的 n，\(V_{kn}\) 已对所有 \(k\in N\) 定义，则由于 \(f_{n,n+1}\) 一致连续，我们可以对 k 用归纳法定义一致邻域 \(V_{k,n+1}\)，使 (3) 成立且

\[
\mathbf {\dot {V}} _ {k + 1, n + 1} ^ {2} \subset \mathbf {V} _ {k, n + 1} \cap \mathbf {U} _ {k + 1, n + 1}
\]

断言得证。

现在设 \(x_0 \in Y_0\)。我们将证明，对每个整数 \(k > 0\)，存在一点 \(z \in X\)，使 \([x_0, f_{\alpha_0}(z)] \in V_{k-1,0}\)；这就证明了定理。由于 \(f_{n,n+1}(Y_{n+1})\) 在 \(Y_n\) 中稠密，我们可以归纳地定义点列 \(x_n \in Y_n\)，使得

\[
\left[ x _ {n}, f _ {n, n + 1} \left(x _ {n + 1}\right) \right] \in \mathrm{V} _ {k + n, n}.\tag{4}
\]

由 (3) 推出，若 \(m \leqslant n\)，则

\[
\left[ f _ {m n} (x _ {n}), f _ {m, n + 1} (x _ {n + 1}) \right] \in \mathbf {V} _ {k + n, m}.\tag{5}
\]

由此我们断定，对固定的 \(m\)，序列 \((f_{mn}(x_n))_{n\geqslant m}\) 是 \(\mathbf{X}_{\alpha_m}\) 中的 Cauchy 序列，从而收敛于一点 \(z_m\)；因为由 (5) 用归纳法可得，对每对整数 \(p\geqslant m,q > 0\)，我们有

\[
\left[ f _ {m p} (x _ {p}), f _ {m, p + q} (x _ {p + q}) \right] \in \mathrm{V} _ {k + p + q - 1, m} \circ \mathrm{V} _ {k + p + q - 2, m} \circ \dots \circ \mathrm{V} _ {k + p, m}\tag{6}
\]

而由 (2) 显然 (6) 的右端含于 \(V_{k+p-1,m}\)。令 q 无限增大；我们特别推出，当 m=p=0 时有 \((x_{0},z_{0})\in\mathbf{V}_{k-1,0}\)，因为 \(V_{k-1,0}\) 是闭的。另一方面，由关系 \(z_{m}=\lim_{n\to\infty}f_{mn}(x_{n})\) 以及 \(f_{m,m+1}\) 的连续性，我们推出 \(f_{m,m+1}(z_{m+1})=z_{m}\)（对每个 \(m\geqslant0\)）。对每个 \(\gamma\in I\)，至少存在一个整数 n 使 \(\alpha_{n}\geqslant\gamma\)；令 \(z_{\gamma} = f_{\gamma, \alpha_n}(z_n)\)，我们立即验证：\(z_{\gamma}\) 不依赖于 \(n\) 的取值（其中 \(\alpha_n \geqslant \gamma\)），且如此定义的族 \((z_\alpha)_{\alpha \in \mathbf{I}}\) 是 \(z\)（\(\mathbf{X} = \varprojlim \mathbf{X}_\alpha\) 的一个点）。由于 \(f_{\alpha_0}(z) = z_0\)，证明完毕。

推论 1 设 \((\mathbf{X}_{\alpha}, f_{\alpha\beta})\) 是以具有可数共尾子集的有向集 I 为指标集的集合逆系统，并设诸 \(f_{\alpha\beta}\) 都是满射。则若 \(X = \varprojlim X_{\alpha}\)，典范映射 \(f_{\alpha}: X \to X_{\alpha}\) 对每个 \(\alpha \in I\) 都是满射。

给每个 \(X_{\alpha}\) 赋以离散一致结构，并应用定理 1。

推论 2 设 I 是具有可数共尾子集的有向集。设 \((\mathbf{X}_{\alpha}, f_{\alpha \beta})\) 与 \((\mathbf{X}_{\alpha}^{\prime}, f_{\alpha \beta}^{\prime})\) 是以 I 为指标集的两个集合逆系统，且对每个 \(\alpha \in I\)，设 \(u_{\alpha}: \mathbf{X}_{\alpha} \to \mathbf{X}_{\alpha}^{\prime}\) 是一个映射，使得诸 \(u_{\alpha}\) 构成映射的逆系统。令 \(u = \varprojlim u_{\alpha}\)。设 \(x = (x_{\alpha}^{\prime})\) 是

\[
\mathrm{X} ^ {\prime} = \varprojlim \mathrm{X} _ {\alpha} ^ {\prime}
\]

的一个元素，它满足下述条件：对每个 \(\alpha \in I\)，存在指标 \(\beta \geqslant \alpha\)，使得对所有 \(\gamma \geqslant \beta\) 有 \(f_{\alpha \gamma}(\overline{u}_{\gamma}^{1}(x_{\gamma})) = f_{\alpha \beta}(\overline{u}_{\beta}^{1}(x_{\beta}^{\prime}))\)。则存在元素 \(x \in X\)，使 \(u(x) = x'\)。

对集合逆系统 \(\overline{u}_{\alpha}^{1}(x_{\alpha}^{\prime})\)（其中每个集合都赋以离散一致结构）应用定理 1。

* 例 设在 C 中给定：(i) 互异点组成的序列 \((a_{n})\)，使序列 \((|a_{n}|)\) 递增且趋于 \(+\infty\)；(ii) 对每个 \(n\)，一个有理函数 \(z \to \mathbb{R}_{n}(z)\)，它定义在 \(\mathbf{C} - \{a_{n}\}\) 中且在 \(a_{n}\) 处有极点；(iii) 以 0 为中心的递增开圆盘序列 \((\mathbf{B}_{n})\)，其并为 \(\mathbf{C}\)，且没有任何 \(a_{k}\) 位于诸圆盘 \(\mathbf{B}_{n}\) 中任何一个的边界上。对每个 \(n\)，令 \(\mathbf{B}_{n}^{\prime}\) 表示 \(\overline{\mathbf{B}}_{n}\) 与 \(\mathbf{C}\) 中由点 \(a_{n}\) 组成的集合的余集之交；并令 \(\mathbf{X}_{n}\) 表示所有映射

\[
z \rightarrow \mathrm{S} (z) = \mathrm{P} (z) + \sum_ {a _ {k} \in \mathrm{B} _ {n}} \mathrm{R} _ {k} (z)
\]

（自 \(\mathbf{B}_n'\) 到 \(\mathbf{C}\)）的集合，其中 \(\mathbf{P}\) 是一个函数在 \(\mathbf{B}_n'\) 上的限制，该函数在 \(\overline{\mathbf{B}}_n\) 上连续且在 \(\mathbf{B}_n\) 内全纯。我们在 \(\mathbf{X}_n\) 中定义度量：

\[
d _ {n} \left(\mathrm{S} _ {1}, \mathrm{S} _ {2}\right) = \sup _ {z \in \mathrm{B} _ {n} ^ {\prime}} | \mathrm{S} _ {1} (z) - \mathrm{S} _ {2} (z) |.
\]

容易验证 \(\mathbf{X}_n\) 关于这个度量是完备的。最后，对 \(n \leqslant m\)，我们定义映射 \(f_{nm}: \mathbf{X}_m \to \mathbf{X}_n\)，使得若 \(\mathbf{S} \in \mathbf{X}_m\)，则 \(f_{nm}(\mathbf{S})\) 是 \(\mathbf{S}\) 在 \(\mathbf{B}_n'\) 上的限制。显然诸 \(f_{nm}\) 都一致连续，且 \((\mathbf{X}_n, f_{nm})\) 是一致空间的逆系统。于是，逆极限 \(X = \varprojlim X_{n}\) 的元素可以典范地恒同于 C 上的一个亚纯函数 F，其仅有的极点是诸点 \(a_{n}\)，并且对每个 n，\(\mathrm{F}(z) - \mathrm{R}_{n}(z)\) 在 \(a_{n}\) 处全纯。经典的 Mittag-Leffler 定理断言 X 非空；由定理 1，我们只需对所有 n 验证条件 \((\mathrm{ML}_{nn})\)。~设

\[
\mathrm{S} _ {n} = \mathrm{P} _ {n} + \sum_ {a _ {k} \in \mathrm{B} _ {n}} \mathrm{R} _ {k}
\]

是 \(\mathbf{X}_n\) 的一个元素，其中 \(\mathbf{P}_n\) 在 \(\overline{\mathbf{B}}_n\) 上连续且在 \(\mathbf{B}_n\) 内全纯；对任意 \(m \geqslant n\)，令 \(\mathbf{Q}_{mn}\) 为

\[
\sum_ {a _ {h} \in \mathrm{B} _ {m} - \mathrm{B} _ {n}} \mathrm{R} _ {h} \quad \text { to } \quad \mathrm{B} _ {n} ^ {\prime};
\]

后一和式是 \(\bar{\mathbf{B}}_n\) 的某个邻域中的全纯函数，因此（由 Taylor 定理）对每个 \(\varepsilon > 0\)，存在多项式 \(\mathbf{P}_{mn}\)，使 \(|\mathbf{Q}_{mn}(\mathbf{Z}) - \mathbf{P}_{mn}(z)| \leqslant \varepsilon\) 在 \(\mathbf{B}_n\) 中成立；若 \(\mathbf{S}_m\) 是 \(\mathbf{S}_n + \mathbf{Q}_{mn} - \mathbf{P}_{mn}\) 在 \(\mathbf{B}_m'\) 上的限制，则我们有 \(\mathbf{S}_m \in \mathbf{X}_m\)，且 \(|\mathbf{S}_m(z) - \mathbf{S}_n(z)| \leqslant \varepsilon\) 在 \(\mathbf{B}_n'\) 中成立。证明完毕。 *

